\documentclass[12pt]{article}

%% Basic document formatting
\usepackage{amsmath, amsthm, amssymb, amsfonts}
\usepackage{mathtools}
\usepackage{xspace}
\usepackage{thmtools}
\usepackage{graphicx}
\usepackage{setspace}
\usepackage{fancyhdr}
\usepackage{titling}
\usepackage[left=0.4in,right=0.4in,top=1in,bottom=1in]{geometry}
\usepackage{float}
\usepackage{hyperref}
\usepackage{mathrsfs}
\usepackage{tabularx}
\usepackage[utf8]{inputenc}
\usepackage[english]{babel}
\usepackage{framed}
\usepackage[dvipsnames]{xcolor}
\usepackage{environ}
\usepackage{tcolorbox}
\tcbuselibrary{theorems,skins,breakable}

\usepackage{enumitem}
\setlist[enumerate]{leftmargin=*}
\setlist[enumerate,1]{labelindent=\parindent}
\setlist[enumerate,2]{labelindent=0pt}

% Blackboard Bold
\newcommand{\Z}{\mathbb{Z}}                 % integers
\newcommand{\Zpl}{\mathbb{Z}_{+}}           % positive integers
\newcommand{\Znn}{\mathbb{Z}_{\geq 0}}      % nonnegative integers. 
\newcommand{\Q}{\mathbb{Q}}                 % rationals
\newcommand{\Qpl}{\mathbb{Q}_{+}}           % positive Rationals
\newcommand{\R}{\mathbb{R}}                 % reals
\newcommand{\Rpl}{\mathbb{R}_{+}}           % positive Reals
\newcommand{\C}{\mathbb{C}}                 % complex numbers
\newcommand{\F}{\mathbb{F}}                 % field

% Words
\newcommand{\st}{\text{ such that }}
\newcommand{\wrt}{\text{ with respect to }}
\newcommand{\with}{\text{ with }}
\newcommand{\ie}{\text{, i.e. }}
\newcommand*{\ditto}{\text{---}\texttt{"}\text{---}}

% Operators
\newcommand{\abs}[1]{\left|#1\right|}                   % absolute value
\newcommand{\norm}[1]{\abs{\abs{#1}}}
\newcommand{\floor}[1]{\left\lfloor #1 \right\rfloor}   % floor
\newcommand{\ceil}[1]{\left\lceil #1 \right\rceil}      % ceiling
\newcommand{\powset}[1]{\mathscr{P}(#1)}

% Category Theory 
\renewcommand{\hom}{\text{Hom}}
\newcommand{\homspace}[3]{\hom_{#1}(#2, #3)}
\newcommand{\coproduct}[9]{
  \begin{tikzcd}[ampersand replacement=\&]
      \& \& #1 \arrow[dll, bend right, "#6"] \arrow[dl, "#5"] \\
   #4 \& #3 \arrow[l, "#9"] \\
      \& \& #2 \arrow[ull, bend left, "#8"] \arrow[ul, "#7"]
  \end{tikzcd}
}

\newcommand{\product}[9]{ 
  \begin{tikzcd}[ampersand replacement=\&]
      \& \& #3 \\
      #1 \arrow[r, "#7"] \arrow[urr, bend left, "#5"] \arrow[drr, bend right, "#6"] \& #2 \arrow[ur, "#8"] \arrow[dr, "#9"] \\ 
      \& \& #4
  \end{tikzcd}
}


% Algebra 
\newcommand{\Syl}[2]{\text{Syl}_{#1}{(#2)}}           % set of Sylow-p subgroups 
\newcommand{\<}{\langle}                % \<x\>, subgroup generated by x 
\renewcommand{\>}{\rangle}
\renewcommand{\index}[2]{[#1 : #2]}
\newcommand{\id}{\text{id}}             % identity element
\newcommand{\lhdneq}{%
  \mathrel{\ooalign{$\lneq$\cr\raise.22ex\hbox{$\lhd$}\cr}}}
\newcommand{\order}[1]{\text{o}(#1)}    % order of an element
\newcommand{\spec}{\operatorname{Spec}}
\newcommand{\rad}{\operatorname{rad}}   % radical of an ideal 
\newcommand{\sym}[1]{\mathfrak{S}_{#1}}
\newcommand{\aut}[1]{\text{Aut}(#1)} 
\newcommand{\gal}[2]{\text{Gal}(#1 / #2)} 
\newcommand{\inn}[1]{\text{Inn}(#1)} 
\let\oldcong\cong
\let\oldequiv\equiv
\renewcommand{\cong}{\oldequiv}
\renewcommand{\equiv}{\oldcong}

\newcommand{\triangleleftneq}{\mathrel{\ooalign{%
  $\trianglelefteq$\cr
  \hidewidth\raise0.06ex\hbox{$\not\mkern4mu$}\hidewidth\cr
}}}

\makeatletter % cycle
\newcommand{\cyc}[1]{(\mathbf{\cyc@process#1\relax})}
\def\cyc@process#1#2\relax{%
	#1%
	\ifx\relax#2\relax
	\else
		\,\cyc@process#2\relax
	\fi
}
\makeatother

\newcommand{\GL}[2]{\text{GL}_{#1}(#2)}                             % general linear group
\newcommand{\SL}[2]{\text{SL}_{#1}(#2)}                             % special linear group
\newcommand{\gl}[2]{\mathfrak{gl}_{#1}(#2)}
\renewcommand{\sl}[2]{\mathfrak{sl}_{#1}(#2)}
\newcommand{\so}[2]{\mathfrak{so}_{#1}(#2)}
\newcommand{\su}[1]{\mathfrak{su}(#1)}

% Linear Algebra
\newcommand{\bmat}[1]{\begin{bmatrix}#1\end{bmatrix}}   % bracketed matrix
\newcommand{\rank}{\operatorname{rank}}                 % rank
\newcommand{\nullity}{\operatorname{nullity}}           % nullity
\newcommand{\tr}{\operatorname{tr}}

% Combinatorics
\newcommand{\qnum}[1]{\left[#1\right]_q}                % q-analog
\newcommand{\qfact}[1]{\left[#1\right]!_q}              % q-analog factorial 
\newcommand{\qbinom}[2]{\binom{#1}{#2}_q}               % Gaussian binomial coefficient

% Topology/Analysis
\newcommand{\ball}[2]{\text{B}_{#1}(#2)}  % B_r(x): open r-balls around x
\newcommand{\diam}{\text{diam}}           % diamter of a set in metric space 
\newcommand{\lowint}[2]{\underline{\int_{#1}^{#2}}}
\newcommand{\upint}[2]{\overline{\int}_{#1}^{#2}}

% Misc Notation
\newcommand{\oneton}[1]{\{1, \ldots, #1\}}
\newcommand{\defeq}{\vcentcolon=}   % :=
\newcommand{\eqdef}{=\vcentcolon}   % =: 
\renewcommand{\bf}[1]{\textbf{#1}}
\newcommand{\im}{\operatorname{im}}
\newcommand{\fnrest}[2]{\left.#1\right|_{#2}}
\newcommand{\isom}{\overset{\sim}{\rightarrow}} 


\renewcommand{\thesection}{\arabic{section}.}
\renewcommand{\thesubsection}{(\alph{subsection})}

\newtheorem{claim}{Claim}
\newtheorem{theorem}{Theorem}
\newtheorem{proposition}{Proposition}
\newtheorem*{lemma}{Lemma}

%% Headers & title setup
\newcommand{\course}{Math140C}
\newcommand{\myname}{Jay Ser}
\setlength{\headheight}{14.5pt}
\pagestyle{fancy}
\fancyhf{}
\renewcommand{\headrulewidth}{0.4pt}
\lhead{\course}
\rhead{\myname}
\cfoot{\thepage}
\setlength{\droptitle}{-4em} 
\title{\course\ - HW \#6}
\author{\myname}
\date{2026.05.31} 

\begin{document}
\maketitle
\thispagestyle{fancy}

%------------------------------------------------------------------------------%
\newcommand{\M}{\mathcal{M}}
\newcommand{\scrC}{\mathscr{C}}
\newcommand{\sigalg}[1]{\sigma\text{-alg}(#1)}

\section*{Q1} 
Let 
\begin{gather*}
  A \defeq \{x \in E \mid f(x) > 0\} \\ 
  E_n \defeq \{x \in E \mid f(x) > 1 / n \}, \quad (n \in \Zpl).
\end{gather*}
Then 
$$A = \bigcup_{n = 1}^{\infty} E_n.$$
Since the measure $\mu$ is monotone and countably additive, 
\begin{equation} \label{eq:1_bidir} 
\mu(A) = 0 \iff \forall n \in \Zpl, \: \mu(E_n) = 0.
\end{equation}
If $f(x)$ is not zero almost everywhere on $E$, then $\mu(A) > 0$. 
By \eqref{eq:1_bidir}, there is some $n \in \Zpl$ such that $\mu(E_n) > 0$. 
By definition of $E_n$, 
$$\frac{1}{n} \chi_n < f$$ 
on $E$. 
Hence 
\begin{align*} 
  \frac{1}{n}\mu(E_n) & = \int_{E}\frac{1}{n} \chi_n d\mu \\  
					  & \leq \int_{E} f d\mu,
\end{align*}
which means $\int_{E} f d\mu > 0$. 
This contradicts the assumption that $\int_{E} f d\mu = 0$, so it must be that $\mu(A) = 0$; 
that is, $f(x) = 0$ almost everywhere on $E$. 

\section*{Q2} 
$E$ is a measurable subset of $E$, so if $f \geq 0$, then $f(x) = 0$ almost everywhere on $E$ by Q1. 
Suppose there is some $x \in E$ such that $f(x) < 0$. 
In other words, letting 
$$A^{nn} = \{x \in E \mid f(x) \leq 0\},$$
$A^{np} \neq \emptyset$. 
$f$ is a measurable function, so $A^{np}$ is a measurable set. 
By assumption on the measurable subsets of $E$, 
\begin{equation} \label{eq:2_int0} 
  0 = \int_{A^{np}} fd\mu = \int_{A^{np}} f^{+} d\mu - \int_{A^{np}} f^{-} d\mu 
\end{equation}
$f^{+} = 0$ on $A^{np}$, so \eqref{eq:2_int0} gives $\int_{A^{np}} f^{-} d\mu = 0$. 
Of course, $f^{-}$ is nonnegative on $A^{np}$, so by Q1, 
\begin{equation} \label{eq:2_0ae} 
  f^{-}(x) = 0 \text{ almost everywhere on } A^{np}.
\end{equation}

Let $A^{-} = \{x \in E \mid f(x) < 0\}$. 
Clearly, $A^{-} \subset A^{np}$. 
Hence the conclusion \eqref{eq:2_0ae} implies $\mu(A^{-}) = 0$. 
Symmetric argument gives $\mu(A^{+}) = 0$ where $A^{+} \defeq \{x \in E \mid f(x) > 0\}$. 
It follows that $f(x) = 0$ almost everywhere on $E$. 


\section*{Q3} 
Let 
\begin{gather*}
  h(x) = \liminf_{n \rightarrow \infty} f_n(x) \\ 
  g(x) = \limsup_{n \rightarrow \infty} f_n(x)
\end{gather*}
Since each $f_n$ is a measurable function, $h$ and $g$ are measurable functions. 
Recall that a sequence has a limit if and only if its liminf and limsup are equal. 
Hence for a given point $x \in X$,
\begin{equation} \label{eq:3_conv} 
  (f_n(x))_{n \in \Zpl} \text{ converges} \iff h(x) = g(x) \iff h(x) - g(x) = 0 
\end{equation} 
Let $k = h - g$. 
$k$ is a measurable function. 
By \eqref{eq:3_conv}, 
\begin{align}
  \{x \in X \mid (f_n(x))_{n \in \Zpl} \text{ converges}\} 
  & = \{x \in X \mid h(x) = 0\} \nonumber \nonumber \\ 
  & = \{x \in X \mid h(x) \geq 0\} \cap \{x \in X \mid h(x) \leq 0\} \label{eq:3_meas}
\end{align}
and \eqref{eq:3_meas} is measurable;
the set of points such that $(f_n(x))$ converges is measurable. 

\section*{Q6} 
It is clear that $f_n \rightarrow 0$ uniformly: 
fix an $\epsilon > 0$ and pick $N \in \Zpl$ such that $1 / N < \epsilon$. 
Then for all $n \geq N$ and for all $x \in \R$, 
$$\abs{f_n(x)} \leq 1 / N < \epsilon.$$ 

However, for every $n \in \Zpl$, 
\begin{align*}
  \int_{-\infty}^{\infty} f_n dx  
  & = \int_{-n}^{n} f_n dx \\ 
  & = \int_{-n}^{n} \frac{1}{n} dx = 2,
\end{align*}
so Lebesgue's dominated convergence theorem is not guaranteed for uniformly convergent sequences of functions even though uniformly convergent sequences are uniformly bounded, i.e., dominated by a single constant function. 

\section*{Q8} 
$$f \in \mathscr{R}_{a}^{b} \iff f \text{ is continuous almost everywhere on } [a, b],$$
and for points $x \in [a, b]$ where $f$ is continuous, 
$$F'(x) = f(x).$$
So assuming $f \in \mathscr{R}_{a}^{b}$, the containment
$$\{x \in [a, b] \mid F'(x) = f(x)\} \supset \{x \in [a, b] \mid f \text{ is continuous on } x\}$$
shows $F'(x) = f(x)$ almost everywhere on $[a, b]$. 



\end{document}
